\documentclass[12pt,a4paper]{article}
\usepackage{amsmath, amssymb, amsthm}
\usepackage[utf8]{inputenc}
\usepackage[T1]{fontenc}
\usepackage{geometry}
\usepackage{enumitem}
\usepackage{xcolor}
\geometry{margin=2.5cm}

\newcommand{\easy}{{\color{green!60!black}$\bigstar$}}
\newcommand{\medium}{{\color{orange}$\bigstar\bigstar$}}
\newcommand{\hard}{{\color{red}$\bigstar\bigstar\bigstar$}}
\newcommand{\pd}[2]{\dfrac{\partial #1}{\partial #2}}

\title{\textbf{Additional Exercises 4}\\
\large Gradient, Directional Derivatives \& Higher-Order Partial Derivatives\\[4pt]
\normalsize Mirrors: List 5\quad|\quad
\easy\ Easy \quad \medium\ Medium \quad \hard\ Hard}
\date{}

\begin{document}
\maketitle
\thispagestyle{empty}

%----------------------------------------------------------
\section*{Exercise 1: Calculate the directional derivative}
%----------------------------------------------------------

\emph{Recall: $D_{\hat{u}}f = \nabla f\cdot\hat{u}$ where $\hat{u}$ is the unit vector.}

\begin{enumerate}[leftmargin=*, label=(\alph*)]

\item \easy\quad
$f(x,y) = x^2+xy-y^2$ at $(1,2)$ in direction $\bar{u}=(3,4)$.

\item \easy\quad
$f(x,y,z) = xy+yz+zx$ at $A(2,1,3)$ in direction $\bar{a}=(3,4,12)$.

\item \easy\quad
$f(x,y) = e^x\cos y$ at the origin in direction making angle $\pi/4$ with the $x$-axis.

\item \medium\quad
$f(x,y,z) = \ln(x^2+y^2+z^2)$ at $A(1,2,1)$ in direction $\bar{a}=(2,4,4)$.

\item \medium\quad
$f(x,y,z) = xyz$ at $A(2,4,2)$ in direction $\bar{u}=(4,2,-4)$.

\item \medium\quad
Find the directional derivative of $f(x,y) = x^2 y$ at $P(1,1)$ in the direction
of the unit vector making angle $\theta = \pi/3$ with the positive $x$-axis.

\item \hard\quad
$f(x,y)=\begin{cases}
\dfrac{x^4 y^2}{x^4+y^2}, & (x,y)\ne(0,0),\\[6pt]
0, & (x,y)=(0,0).
\end{cases}$
at $(0,0)$ in direction $\bar{u}=(1,1)$.\\
\emph{(Use the limit definition --- the formula $\nabla f\cdot\hat{u}$ may not apply here.)}

\end{enumerate}

%----------------------------------------------------------
\section*{Exercise 2: Find the direction (if it exists) for the given rate of change}
%----------------------------------------------------------

\begin{enumerate}[leftmargin=*, label=(\alph*)]

\item \easy\quad
$f(x,y) = 3x + 4y$ at any point. Find the direction of steepest ascent and its rate.

\item \medium\quad
$f(x,y) = x^2 y + y^3$ at $(1,-1)$.
Find a unit vector $\hat{u}$ such that $D_{\hat{u}}f = 0$.

\item \medium\quad
$f(x,y,z) = xy+z$ at $(0,1,-2)$.
Find $\hat{u}$ such that $D_{\hat{u}}f = 1$.
Does a direction with $D_{\hat{u}}f = 20$ exist?

\item \hard\quad
$f(x,y) = x^2 y + y^3$ at $(1,-1)$.
Find $\hat{u}$ such that $D_{\hat{u}}f = 20$, or show no such direction exists.

\end{enumerate}

%----------------------------------------------------------
\section*{Exercise 3: Find the gradient of the scalar field at point $M_0$}
%----------------------------------------------------------

\begin{enumerate}[leftmargin=*, label=(\alph*)]
\item \easy\quad $u = x+y+2xy$ at $M_0(1,1)$
\item \easy\quad $u = x^2+y^2+z^2$ at $M_0(2,0,0)$
\item \medium\quad $u = \ln(x^2+y^2+z^2)$ at $M_0(1,1,1)$
\item \medium\quad $u = x^2 e^{y-z}$ at $M_0(2,1,1)$
\item \medium\quad $u = \arctan\!\left(\dfrac{y}{x}\right)+\arctan\!\left(\dfrac{x}{y}\right)$.
Compute $\nabla u$ and explain the result.
\item \hard\quad $u = \arctan(xyz)+\operatorname{arccot}(xyz)$.
Compute $\nabla u$ and explain the result.
\item \hard\quad $u = f(r)$ where $r = \sqrt{x^2+y^2+z^2}$, $f$ differentiable.
Express $\nabla u$ in terms of $f'(r)$ and $\bar{r} = (x,y,z)$.
\end{enumerate}

%----------------------------------------------------------
\section*{Exercise 4: Second-order partial derivatives}
%----------------------------------------------------------

\emph{Verify Schwarz's theorem where asked: $f_{xy}=f_{yx}$.}

\begin{enumerate}[leftmargin=*, label=(\alph*)]
\item \easy\quad $z=\sin(xy)$. Find $\dfrac{\partial^2 z}{\partial x\,\partial y^2}$.
\item \easy\quad $z = x^3 y + x y^3$. Verify $f_{xy} = f_{yx}$.
\item \medium\quad $z=x^y$. Verify $f_{xy}=f_{yx}$.
\item \medium\quad $u=\arctan(y/x)$. Verify the Laplace equation:
$\dfrac{\partial^2 u}{\partial x^2}+\dfrac{\partial^2 u}{\partial y^2}=0$.
\item \medium\quad $u = e^x(x\cos y - y\sin y)$. Verify $u_{xx}+u_{yy}=0$.
\item \hard\quad $u=\ln\dfrac{1}{r}$, $r=\sqrt{(x-a)^2+(y-b)^2}$.
Verify that $u$ satisfies the Laplace equation for $(x,y)\ne(a,b)$.
\end{enumerate}

%----------------------------------------------------------
\section*{Exercise 5: Find the second-order differential $d^2 u$}
%----------------------------------------------------------

\emph{Recall: $d^2 u = u_{xx}(dx)^2+2u_{xy}\,dx\,dy+u_{yy}(dy)^2$.}

\begin{enumerate}[leftmargin=*, label=(\alph*)]
\item \easy\quad $u = x^2 y^3$
\item \easy\quad $u = e^{xy}$
\item \medium\quad $u = x^y$
\item \medium\quad $u = x\ln(xy)$
\item \medium\quad $u = \sin(x+y) + \cos(x-y)$
\item \hard\quad $u = f(x^2-y^2)$, $f$ twice-differentiable.
Express $d^2 u$ in terms of $f'$ and $f''$.
\item \hard\quad $u = \ln\!\sqrt{x^2+y^2}$. Compute $d^2 u$ and verify $u_{xx}+u_{yy}=0$.
\end{enumerate}

\bigskip
\begin{center}
\small\itshape
\easy\ Direct formula\quad
\medium\ Gradient properties or Schwarz\quad
\hard\ Definition from first principles or abstract function.
\end{center}

\end{document}
